\documentclass[11pt]{article}
\usepackage[a4paper,margin=27mm]{geometry}
\usepackage{amsmath,amssymb,mathtools}
\usepackage{fontspec,unicode-math}
\setmainfont{Latin Modern Roman}
\setmathfont{Latin Modern Math}
\usepackage{xurl}
\usepackage[hidelinks]{hyperref}
\newcommand{\sref}[2]{\hyperlink{source:#1:#2}{[S#2]}}
\newcommand{\cataloguescope}[1]{\paragraph{Mathematical question.}#1}
\hypersetup{pdftitle={Problem statement}}
\begin{document}
\section{Are there infinitely many prime quadruplets}\label{statement}

\subsection{Definitions and mathematical statement}

\paragraph{Shared notation.}$\mathbb N=\{1,2,\ldots\}$, and $\mathbb Z,\mathbb Q,\mathbb R,\mathbb C$ denote the integers, rationals, reals and complex numbers. Any other conventions are specified locally.


A prime quadruplet is a tuple $(p,p+2,p+6,p+8)$ whose four entries are prime.
\paragraph{Statement recorded in the supplied source.}
Are there infinitely many primes $p$ for which all four entries of this tuple are prime? \sref{statement}{2}
\subsection{Short English statement}
Do infinitely many primes start the four-prime pattern with offsets zero, two, six and eight?
\subsection{Sources}
\begin{description}
\item[\hypertarget{source:statement:1}{S1}] ULAM.AI and the UnsolvedMath Contributors, UnsolvedMath: A Curated Collection of Open Mathematics Problems, catalogue entry. CC BY 4.0. \url{https://huggingface.co/datasets/ulamai/UnsolvedMath}.
\item[\hypertarget{source:statement:2}{S2}] Wikipedia contributors, \emph{List of unsolved problems in mathematics}, revision 1366636767 (2026) Number theory, Are there infinitely many prime quadruplets. \url{https://en.wikipedia.org/w/index.php?title=List_of_unsolved_problems_in_mathematics&oldid=1366636767#Number_theory}.
\item[\hypertarget{source:statement:3}{S3}] OEIS Foundation and contributors, \emph{The On-Line Encyclopedia of Integer Sequences}, A007530: Prime quadruples: numbers k such that k, k+2, k+6, k+8 are all prime., official data snapshot \texttt{693b5499c9d2} (2026) Name, definition and comments. \url{https://raw.githubusercontent.com/oeis/oeisdata/693b5499c9d248f570dd2d979e8bd1d7eecf689b/seq/A007/A007530.seq}.
\end{description}
\label{end:statement}
\end{document}
